The torsion condition is the equation $D\wedge\hb^{T}(f) = \hb^{T}(\nabla\wedge f)$ of Definition \claimref{obj.synthesis.gtg-flat-spacetime-gauge-fields}(f); for the background vector fields $f^a$ with $\hb^{T}(f^a) = \gamma^a$, a constant physical frame, $D_{\hb(\gamma_b)}\gamma^a = \omega(\gamma_b)\cdot\gamma^a$ gives its left side as $\sum_b \gamma^b\wedge(\omega(\gamma_b)\cdot\gamma^a)$, and $\Omega(a) = \sum_c (\gamma^c\cdot\hb^{-1}(a))\,\omega(\gamma_c)$ on background directions (Fact \claimref{synthesis.imported-gtg-dictionary}(i) and (ii); the one-forms $\omega^a{}_b$ of the coframe are the components of this equation, a labelled reading). (1) $\beta = (2M)^{1/2} r^{-1/2}$ gives $\beta' = -\tfrac12 (2M)^{1/2} r^{-3/2} = -\beta/2r$ and $\beta\beta' = -\beta^2/2r = -M/r^2$. (2) The physical frame of $\hb = 1 - \beta P$ is $\hb^{-1}(e_t) = \gamma_0 + \beta\gamma_1$, $\hb^{-1}(e_r) = \gamma_1$, $\hb^{-1}(e_\vartheta) = r\gamma_2$, $\hb^{-1}(e_\varphi) = r\sin\vartheta\,\gamma_3$, so $f^0 = \nabla t$, $f^1 = \beta\nabla t + \nabla r$, $f^2 = r\nabla\vartheta$, $f^3 = r\sin\vartheta\,\nabla\varphi$, and $\hb^{T}(\nabla r) = \gamma^1 - \beta\gamma^0$, which is where the gauge field's off-diagonal character enters the connection. On the left, $\nabla\wedge f^1 = \beta'\,\nabla r\wedge\nabla t$ and $\nabla\wedge f^2 = \nabla r\wedge\nabla\vartheta$, whose images under the outermorphism $\hb^{T}$ are $\beta'\,\gamma^1\wedge\gamma^0$ and $\tfrac1r(\gamma^1 - \beta\gamma^0)\wedge\gamma^2$. On the right, with $(\gamma_1\wedge\gamma_0)\cdot\gamma^1 = -\gamma_0$, the term $\omega(\gamma_1) = -\beta'\sigma_r$, $\sigma_r = \gamma_1\gamma_0$, gives $\gamma^1\wedge\beta'\gamma_0 = \beta'\,\gamma^1\wedge\gamma^0$ for $a = 1$; the term $\tfrac1r(\beta\gamma_0 + \gamma_1)\wedge\gamma_2$ of $\omega(\gamma_2)$ gives $\gamma^2\wedge\tfrac1r(\beta\gamma_0 + \gamma_1) = \tfrac1r(\gamma^1 - \beta\gamma^0)\wedge\gamma^2$ for $a = 2$ and multiples of $\gamma^2\wedge\gamma_2 = 0$ for $a = 0, 1$; and $\omega(\gamma_0) = 0$ with $\nabla\wedge f^0 = 0$ for $a = 0$. The $\gamma_3$ leg is the same computation with $f^3$ and adds the round sphere's term $\tfrac{\cot\vartheta}{r}\gamma_2\wedge\gamma_3$. The script verifies the stated $\Omega$ against the torsion condition exactly at six points, with uniqueness by the rank of the torsion map, and the $\theta^0$, $\theta^1$ and $\theta^2$ components as exact identities in $\mathbb{Q}(r)[\beta]$. (3) $\Omega(e_t) = \omega(\gamma_0) + \beta\,\omega(\gamma_1) = -\beta\beta'\,\sigma_r$ and $\Omega(e_r) = \omega(\gamma_1) = -\beta'\,\sigma_r$, and the numerical identities are (1). (In components, a labelled reading: the frame duals of \eqref{eq:pgcoframe} satisfy $\partial_t = \hat e_0 + \beta\,\hat e_1$ and $\partial_r = \hat e_1$, so $\omega^1{}_0(\partial_t) = -\beta'\beta$ and $\omega^1{}_0(\partial_r) = -\beta'$.)
